\section{Keeping Faith with the Dying}

A dying person remains our neighbor, entitled to the patience and reverence charity gives to a person. The changing body can dominate the room: breathing, discomfort, medication, and the labor of ordinary care demand attention. Within those necessities there is still someone to address by name, someone whose preferences matter, whose history deserves respect, and whose weakness does not cancel dignity. The body needs care precisely because it belongs to this person, not because it can still perform the actions that once made the person useful to others.

Accompaniment often asks us to accept a slower pace than our anxieties would prefer. A conversation may need to stop before we feel it is finished. The dying may want silence when visitors have prepared words, or may return repeatedly to a worry the family believed was already settled. We should listen for the person rather than insist upon a scene that reassures us. A confession of fear may be an act of trust toward the listener. It need not be immediately corrected into a statement of confidence.

The Lord asked his disciples to watch with him in his agony. We can remember his request when remaining present seems a small service beside the suffering we cannot relieve. Watching may include prayer, but also water, a changed position, help with a practical concern, or relief for another exhausted watcher. Caregivers need care as well. The relative who has remained attentive for weeks can become weary, irritable, and ashamed of weariness. Taking a turn at the bedside or helping with other duties allows that person to rest without leaving the dying alone.

Christian care protects life without treating every possible prolongation as obligatory. Ordinary care remains owed even when death is imminent. Medical procedures that are burdensome, dangerous, extraordinary, or disproportionate to their expected outcome may legitimately be discontinued: one accepts the inability to prevent death without choosing to cause it. Pain relief may also be legitimate when a risk of shortening life is foreseen, provided death is not intended as either means or end. The Catechism distinguishes these decisions from euthanasia and assigns them to the competent patient or those legitimately entitled to act, respecting the patient's reasonable will and interests (CCC 2276--2279). Applying the principles requires the actual prognosis, burdens, alternatives, and competent medical and pastoral counsel. Inability to cure, by itself, does not make all treatment or care dispensable.

These distinctions serve a person whom love must neither abandon nor claim as its possession. A family may find it hard to accept that further intervention cannot secure the return it longs for. Another may be tempted to treat dependence as the disappearance of meaningful life. In both situations, attentive care remains a duty. We are called to seek the person's good within the truth of the condition, while acknowledging that our power to heal has limits. Hope in resurrection permits us to recognize those limits without making death our chosen solution.

The dying person may still be able to offer gratitude, forgive, ask pardon, or entrust a concern. These possibilities deserve room, but they should not become demands for a visibly edifying death. Pain, delirium, exhaustion, or unconsciousness can take away the power to speak coherently. The final words we hear are not a complete judgment upon a lifetime. We continue to pray, seek the Church's ministry where possible, and commend the person to the God who knows the heart more deeply than any witness beside the bed.

Even a solitary or sudden death is not outside that knowledge. Those who were absent may suffer over what they were unable to say or do. There may be real faults to acknowledge, but absence itself need not be turned into a claim that the loved one was therefore abandoned by God. The Lord's presence does not depend upon our success in being present. We can entrust the unspoken words, repent of what needs repentance, and continue the charity that remains possible through prayer and care for those who survive.

